\section{Pattern counts and a nearly uniform outer coloring}
\label{sec:outer}

\subsection{Counting comparisons, with equality retained}

\begin{lemma}[Arrangement count]\label{lem:arrangement}
The joint signs, including zero signs, of \(u\) affine functions on
\(\R^e\) have at most
\begin{equation}\label{eq:arrangement}
 \sum_{d=0}^e\binom ud\sum_{j=0}^{e-d}\binom uj
 \le(e+1)^2(u+1)^e
\end{equation}
realized values. Restricting the parameter domain does not increase
this count.
\end{lemma}

\begin{proof}
The strict regions of \(u\) hyperplanes in dimension \(e\) number at
most \(\sum_{j=0}^e\binom uj\). Add the hyperplanes one at a time:
a new hyperplane splits no more regions than the number of regions
in its induced arrangement of dimension \(e-1\). Induction and
Pascal's identity prove the bound.
For a sign pattern with zero signs, select \(d\) independent zero
equations defining their common affine intersection. There are at
most \(\binom ud\) choices. On that intersection count strict regions
in dimension \(e-d\); all equations vanishing identically there have
forced zero signs. Sum over \(d\). The final estimate follows by
bounding every binomial product by \((u+1)^e\).
\end{proof}

This elementary form of the standard arrangement bound
\cite{stanley2007} keeps exact boundary events visible.

\begin{lemma}[Global label entropy]\label{lem:global}
Let \(T_{\rm glob}\) be the number of literal label words
\((L(n_j))_{j<k}\) of nonzero-step cyclic progressions. Then
\begin{equation}\label{eq:global-entropy}
 \log T_{\rm glob}
 =O\!\left(D(\log k+\log\lambda)\right)
 =O_A(k^{2/3}L^{5/3}).
\end{equation}
\end{lemma}

\begin{proof}
For a circle partition with \(F\) endpoints, freely vary a scalar
start and step \(a,b\in[0,1)\). The labels of
\((a+jb)\bmod1\), \(j<k\), are determined by comparisons of \(a+jb\)
with all lifts \(\tau+z\) of endpoints, where \(-1\le z\le k\).
At most \(3k^2F\) affine lines in the parameters \((a,b)\) suffice.
By \cref{lem:arrangement}, there are at most
\(9(3k^2F+1)^2\) scalar words.

Apply this separately to all \(2D\) coordinates. Letting their starts
and steps vary independently only enlarges the possible set.
We may take \(F=|U|+|V|=k^2\lambda+O(k^3L)\), so
\[
 T_{\rm glob}\le[9(3k^2F+1)^2]^{2D}.
\]
Now use \cref{lem:dilation}. The identity of every interval, not just
its rank among the visits, is preserved in this count.
\end{proof}

\subsection{Eligible blocks and counting through one fixed label}

Fix \(h_0<h\le2M\) and \(t\in\{0,\ldots,h-1\}^D\).
In the second coordinate system define
\[
 s_i=\frac{h}{\gcd(h,\lambda t_i)}.
\]
A coordinate is \emph{stationary} if \(s_i=1\), and \emph{rotating}
otherwise. By \cref{lem:dilation}, a rotating coordinate has \(s_i>h_0\).
Stationary refers only to its rational motion; a real drift is allowed.

\begin{definition}\label{def:eligible}
Consider paths, in the torus,
\[
 X_z=x+\frac zh(t+u),\qquad
 Y_z=y+\frac zh(\lambda t+v),\qquad 0\le z<h,
\]
where \(x,y\in\T^D\) and \(u,v\in\R^D\).
A position is \emph{regular} when
\(\rho(Y_{z,i})\ge\eta\) in every rotating coordinate.
Record \(L(X_z,Y_z)\) at a regular position and a dummy symbol at
every other position.
The resulting literal word is \emph{eligible} for \(h,t\) if it has
a realization satisfying all of the following:
\begin{enumerate}
 \item at least \(h/2\) positions are regular, and their full labels
       are pairwise distinct;
 \item \(|u_i|\le H_x\) in every coordinate and \(|v_i|\le H\)
       in each rotating coordinate;
 \item in each stationary coordinate,
       \(|v_i|\le |V(Y_{z,i})|\) at every regular position.
\end{enumerate}
\end{definition}

Only the recorded word and \(h,t\) are counted; its potentially many
real realizations do not give extra events.

\begin{lemma}[Anchored local count]\label{lem:local}
There is an absolute constant \(C_1\) such that for any full label
\(\beta\) and \(h_0<h\le2M\), the number of pairs consisting of \(t\)
and an eligible pattern for \(h,t\) containing \(\beta\) is at most
\begin{equation}\label{eq:local-count}
 \left(\frac{C_1 D h}{\zeta}\right)^{6D}.
\end{equation}
The constant is independent of \(q,\lambda\), and the widths of the
intervals in \(\beta\).
\end{lemma}

\begin{proof}
Fix \(t\) and an anchor index \(z_0\) recording \(\beta\).
In any scalar coordinate re-anchor the representative at \(z_0\).
Its value at \(z\) before reducing modulo one is
\[
 a+\tau_z+\frac{z-z_0}{h}b,
\]
where \(\tau_z\) is fixed, \(a\) belongs to the interval in \(\beta\),
and \(b\) is the appropriate drift.
We retain only the restrictions on \(a,b\) imposed at the anchor.
For a first coordinate the total target range has length at most
\(3H_x\), hence meets at most four lifted \(U\)-endpoints.
For a stationary second coordinate with anchor interval \(I\), the
translation is integral and \(|b|\le |I|\).
Every target lies in \(\mathcal N(I)\), meeting at most \(49\)
endpoints by \cref{lem:mesh}.
For a rotating second coordinate the target range has length at most
\(4H\), and its truncated output is determined by at most
\(C_0D/\zeta\) breakpoints.

All the enclosing ranges and breakpoint lifts are fixed before
\((a,b)\) vary. Their magnitude introduces no additional factor:
each range is shorter than one in the second system, so each
relevant circle breakpoint has at most one lift in it.
Over all \(h\) positions there are \(O(Dh/\zeta)\) comparisons in
two real parameters. Thus the scalar output has at most
\(C_2(Dh/\zeta)^2\) possibilities, including boundary outputs.
Taking the product over \(2D\) scalar coordinates determines the
dummy mask as well as the actual regular labels.
There is no separate mask to enumerate.

There are \(h^D\) choices of \(t\) and at most \(h\) choices of \(z_0\).
The resulting bound is
\[
 h^{D+1}[C_2(Dh/\zeta)^2]^{2D}
 \le(C_1Dh/\zeta)^{6D},
\]
after enlarging an absolute constant.
\end{proof}

\subsection{Residual tests: the change needed for near-uniformity}

For each realized word \(w\) and integer \(h_0<h\le2M\), divide its
first \(h\fl{k/h}\) positions into consecutive blocks of length \(h\).
Let
\begin{equation}\label{eq:residual-set}
 S(w,h)=\bigcup\{\text{full blocks containing no heavy position of }w\}.
\end{equation}
Every position in this set is light. Crucially, \(S(w,h)\) is determined
by \(w,h\), rather than being an arbitrary subset of the light
positions.

We impose two global kinds of tests: the set of all light positions,
and each \(S(w,h)\). A test is included only if its set has at least
\(\zeta k\) positions. For such a set \(S\), let \(w_{\max}\) be its
maximum label multiplicity. Group occurrences by label, fix their
order deterministically from the word, and distribute the list
cyclically into \(w_{\max}\) rows.
Each row has distinct labels and at least
\begin{equation}\label{eq:row-length}
 \fl{|S|/w_{\max}}\ge\zeta M-1
\end{equation}
positions. The number of row tests is at most
\(k(2M+1)T_{\rm glob}\).

The local tests are all eligible recorded patterns, tested on their
regular positions. Each also has distinct labels.
For a test of length \(l\), we require each color \(b\in\Z/s\Z\)
to occur between \((1/s-\zeta)l\) and \((1/s+\zeta)l\).

\begin{lemma}[Finite local lemma]\label{lem:lll}
Let finitely many events \(E\) have a dependency graph such that
each \(E\) is independent of the joint sigma-algebra of its
nonneighbors. If \(0<z_E<1\) and
\[
 \Prb(E)\le z_E\prod_{F\sim E}(1-z_F)
\]
for every \(E\), then all their complements have positive joint
probability.
\end{lemma}

\begin{proof}
An induction on the number of conditioned complements proves
\(\Prb(E\mid\bigcap_{F\in T}F^c)\le z_E\) whenever \(E\notin T\),
and proves that every such conditioning event has positive probability.
For the conditional estimate split \(T\) into neighbors \(T_1\) of
\(E\) and nonneighbors \(T_2\).
If \(T_1\) is empty, use independence.
Otherwise expose the complements in \(T_1\) one at a time after
conditioning on \(T_2\). The induction hypothesis gives a denominator
at least \(\prod_{F\in T_1}(1-z_F)\), while independence bounds the
numerator by \(\Prb(E)\).
The displayed hypothesis completes the induction.
Positivity follows by multiplying the conditional probabilities of
successive complements. This is the finite asymmetric local lemma
\cite{erdosLovasz1975}.
\end{proof}

\begin{proposition}[Simultaneously balanced tests]\label{prop:outer-tests}
For sufficiently large absolute \(A\), and then all sufficiently large
\(k\), there is a coloring \(c_*:\cL\to\Z/s\Z\) satisfying every
global row test and every eligible local test above, for \(s=2\)
or \(s=3\).
\end{proposition}

\begin{proof}
Color each label independently and uniformly in \(\Z/s\Z\).
For a fixed color its count on \(l\) distinct labels is a sum of
independent Bernoulli variables of mean \(1/s\).
For a Bernoulli variable \(X\) with mean \(q\), the function
\[
 g(t)=\log(1-q+qe^t)-qt
\]
has \(g(0)=g'(0)=0\) and \(g''(t)\le1/4\).
Hence \(\E e^{t(X-q)}\le e^{t^2/8}\).
Markov's inequality with \(t=4\zeta\), and the analogous negative
value, gives a two-sided tail at most \(2e^{-2\zeta^2l}\).
Union over the at most three colors bounds each bad event by
\begin{equation}\label{eq:hoeffding}
 \Prb(E)\le6e^{-2\zeta^2l(E)}.
\end{equation}

Assign \(z_E=e^{-\zeta^2l(E)/2}\).
For any one label, the total incident weight from global events is
at most
\begin{equation}\label{eq:global-weight}
 k(2M+1)T_{\rm glob}
 \exp[-\zeta^2(\zeta M-1)/2]=o(\zeta^2).
\end{equation}
Indeed its positive logarithm is \(O_A(k^{2/3}L^{5/3})+O(L)\),
whereas its negative term has order \(k^{3/4}/L^3\).
The latter grows faster.
For local events \cref{lem:local} bounds the incident weight by
\begin{equation}\label{eq:local-weight}
 \sum_{h>h_0}
 (C_1Dh/\zeta)^{6D} e^{-\zeta^2h/4}.
\end{equation}
The function \(\log(C_1Dh/\zeta)/h\) is decreasing for \(h\ge h_0\)
once \(k\) is large. At the initial point,
\[
 \frac{6D\log(C_1Dh_0/\zeta)}{h_0}
 =\left(\frac4A+o(1)\right)\zeta^2.
\]
Choose \(A\) large enough that this is at most \(\zeta^2/8\).
Then \eqref{eq:local-weight} is at most
\[
 \sum_{h>h_0} e^{-\zeta^2h/8}
 =O\!\left(\zeta^{-2}e^{-\zeta^2h_0/8}\right)
 =o(\zeta^2).
\]
These estimates also show \(\zeta^2l(E)\to\infty\) uniformly over the
included tests.

Join events sharing a label. They have the required dependency
property because the underlying label colors are independent.
For large \(k\), the total weight incident to any label is at most
\(\zeta^2/16\), so
\(\sum_{F\sim E}z_F\le\zeta^2l(E)/16\).
All weights are less than \(1/2\), and
\(\log(1-z)\ge-2z\) on that interval. Consequently
\[
 z_E\prod_{F\sim E}(1-z_F)
 \ge e^{-5\zeta^2l(E)/8}
 \ge6e^{-2\zeta^2l(E)}
 \ge\Prb(E).
\]
Apply \cref{lem:lll}, and fix one successful coloring \(c_*\).
\end{proof}

Define \(c_0(n)=c_*(L(n))\).
All randomness used for this outer coloring is now finished.

\begin{theorem}[Nearly uniform or globally affine]\label{thm:dichotomy}
For every nonzero-step cyclic progression \(n_j=n_0+jd\), either
its entire centered second-coordinate path is affine in \(\R^D\),
or, for every color \(b\in\Z/s\Z\),
\begin{align}
 \#\{j:c_0(n_j)=b\}&\ge(1/s-4\zeta)k,\label{eq:outer-each}\\
 \#\{j:c_0(n_j)\ne b\}&\ge(1-1/s-4\zeta)k.\label{eq:outer-complement}
\end{align}
\end{theorem}

\begin{proof}
If all labels are light, the global test on all \(k\) positions proves
both assertions, with \(\zeta\) in place of \(4\zeta\).
Otherwise use a heavy label to obtain \(h,t,u,v\) from
\cref{lem:return}. When \(h\le h_0\), \cref{lem:short} gives the
affine alternative.

Suppose \(h>h_0\). For a rotating coordinate, the rational part of a
full \(h\)-block visits an equally spaced \(s_i\)-point grid, each point
\(h/s_i\) times, where \(s_i>h_0\).
The actual drift has magnitude at most \(4MH/k<\eta\).
Nonregular positions in that coordinate therefore correspond to grid
points within \(2\eta\) of the cut.
An arc of length \(4\eta\) contains at most \(4\eta s_i+1\) such
grid points, including endpoints.
The nonregular fraction in a block is thus at most
\[
 D(4\eta+1/h_0)
 =\zeta/25+D/h_0\le\zeta.
\]
Distinct block indices have distinct full labels by
\cref{lem:return}. Also \(|u_i|\le H_x\) and \(|v_i|\le H\)
eventually.

Any full block containing a heavy position \(j_*\) is eligible.
To check its only remaining condition, take a stationary coordinate,
and put \(I_*=V(y_i(n_{j_*}))\), \(w_*=|I_*|\).
The circle displacement within the block is at most
\(|v_i|\le(2M/k)w_*\).
By the neighborhood estimate of \cref{lem:mesh}, every interval
visited in that block has length at least \(w_*/16\).
Since \(2M/k\le1/16\) eventually, its length is at least \(|v_i|\).
This proves eligibility even if the heavy anchor itself is nonregular.

Let \(B_*\) be the union of heavy-containing full blocks, and let
\(S=S(w,h)\) be their light complement among the full blocks.
The remaining tail has fewer than \(h\le2M\le\zeta k\) positions.
For either the indicator of one color or its complement, let its
uniform mean be \(\beta\), where \(\beta=1/s\) or \(1-1/s\).
The local tests give at least
\((\beta-\zeta)(1-\zeta)|B_*|\) successes on \(B_*\).
If \(|S|\ge\zeta k\), its row tests add
\((\beta-\zeta)|S|\), giving at least
\((\beta-\zeta)(1-\zeta)^2k\).
If \(|S|<\zeta k\), discard \(S\) and the tail instead, to get
\((\beta-\zeta)(1-\zeta)(1-2\zeta)k\).
Both expressions are at least \((\beta-4\zeta)k\) for \(0\le\beta\le1\).
This proves \eqref{eq:outer-each} and \eqref{eq:outer-complement}.
\end{proof}

\begin{remark}
The two tested parts \(B_*\) and \(S(w,h)\) are disjoint.
It would be invalid to replace the test on \(S(w,h)\) by balance on
the entire light set and then add its count to the block counts:
light positions already lying in \(B_*\) could be counted twice.
The new residual tests resolve exactly this overlap.
\end{remark}
